\documentclass[10pt,a4paper]{amsart}
\usepackage[margin=19mm]{geometry}
\usepackage{amsmath,amssymb,mathtools}
\usepackage[scr=rsfs,cal=euler]{mathalfa}
\usepackage{xfrac}
\usepackage[unicode,hidelinks,bookmarks=false]{hyperref}
\newcommand{\ie}{i.e.\ }
\newcommand{\cf}{cf.\ }
\newcommand{\resp}{respectively\ }
\newcommand{\Subsection}{\subsection}
\providecommand{\nobreakdash}{\nobreak\hbox{-}}
\setlength{\emergencystretch}{2em}
\multlinegap0pt
\usepackage{amssymb}
\usepackage{csquotes}
\usepackage{stackengine}
\usepackage{xcolor}
\usepackage{tikz}
\usepackage{enumerate}
\usepackage{mathtools}
\usepackage{breqn}
\newtheorem{proposition}{Proposition}
\newcommand\V{\mathrm{V}_{\Omega}}
\title{Kernel zero-sets, quantum graph ideals, and Hadamard graphons}
\author{Madelyn Andersen}
\date{Source: arXiv:2112.03885v3 (CC BY 4.0)}
\begin{document}
\maketitle

\noindent\emph{Source and licence.} Kernel zero-sets, quantum graph ideals, and Hadamard graphons. Version arXiv:2112.03885v3 (CC BY 4.0), distributed by arXiv under the Creative Commons Attribution 4.0 International licence, http://creativecommons.org/licenses/by/4.0/, as recorded in the arXiv licence metadata for this exact version and shown on the abstract page https://arxiv.org/abs/2112.03885v3. Full licence notice: \texttt{licenses/arxiv-2112.03885-CC-BY-4.0.txt}.\par\smallskip

\noindent\textit{Editorial scope.} Counted unit: the article's proposition topology (source0.tex lines 368--370). The article's complete original proof of it is delivered with it (source0.tex lines 372--402, one \texttt{proof} environment of 31 lines). No mathematics is written, repaired or re-derived: the statement and the proof are the article's own lines, in the article's own notation, and no retained line is substituted or re-flowed.  No further source-local context is needed: every cross-reference of every retained line resolves inside the two delivered spans.  Nothing was omitted from any retained span, so the package carries no omission record.  No retained line contains a citation, so the wrapper carries no bibliography and no bibliography entry.  No retained line contains a figure environment, an image-inclusion command or drawing code, so no image is delivered and the package declares no asset dependency.  Editorial formatting only, in addition: an AMS \texttt{amsart} preamble that restates the article's own declarations and macros used by the retained lines, namely the environments \texttt{proposition} on separate counters, together with the standard packages those lines need.  The retained environments print in this document as Proposition 1 in order of appearance, whereas the article prints the same items with the numbering of its own full document.  The complete native source of the article as distributed in the arXiv e-print for this exact version is bundled unchanged as \texttt{sources/119365/source0.tex}.
\begin{proposition}\label{proposition: zariski topology}
The algebraic subsets of $\Omega$ are closed under finite unions and arbitrary intersections. Moreover, $\varnothing$ and $\Omega$ are algebraic. Hence the complements of algebraic subsets define a topology on $\Omega$, called the Zariski topology associated to homomorphism-density zero-sets.
\end{proposition}

\begin{proof}
Let $Y_1=\V(Q_1)$ and $Y_2=\V(Q_2)$. Define
\[
    Q_1Q_2:=\{gh:g\in Q_1,\ h\in Q_2\}.
\]
We claim that
\[
    Y_1\cup Y_2=\V(Q_1Q_2).
\]
If $W\in Y_1\cup Y_2$, then either all elements of $Q_1$ vanish at $W$, or all elements of $Q_2$ vanish at $W$. Hence every product $gh$ with $g\in Q_1$ and $h\in Q_2$ vanishes at $W$.

Conversely, suppose $W\in\V(Q_1Q_2)$ and $W\notin Y_1$. Then there exists $g\in Q_1$ such that $t(g,W)\neq0$. For every $h\in Q_2$,
\[
    0=t(gh,W)=t(g,W)t(h,W),
\]
so $t(h,W)=0$. Hence $W\in Y_2$. This proves the finite-union claim.

For arbitrary intersections,
\[
    \bigcap_{\alpha}\V(Q_\alpha)=\V\!\left(\bigcup_{\alpha}Q_\alpha\right).
\]
Finally,
\[
    \varnothing=\V(\{K_0\})=\V(\{K_1\}),
\]
because $t(K_0,W)=t(K_1,W)=1$ for every kernel $W$, while
\[
    \Omega=\V(\varnothing)=\V(\{0\}).
\]
Thus algebraic sets satisfy the axioms for closed sets of a topology.
\end{proof}

\end{document}
